\section{Introduction}
\label{introduction}
In recent years, autonomous systems have attracted increasing attention in both academia and industry. 
Light Detection and Ranging (LiDAR) is one of the key sensors in various autonomous systems. 
It provides precise and rich geometric information about the surrounding environment \cite{li2024di, li2023lwsis, SECOND, cheng2023language, yin2024fusion} and has been widely applied across a range of fields \cite{shi2020pv, shi2019pointrcnn,malavazi2018lidar, weiss2011plant, resop2019drone, zhang2014loam}.
Unlike the extensively studied image modality, research on LiDAR is still limited.
Existing studies have focused on two main directions: one is generating high-quality LiDAR data \cite{LiDARGen2022, xiong2023ultralidar, nakashima2024lidar} to reduce the high cost of real-world acquisition \cite{meng2020weakly, wu2020deep}, and the other is predicting LiDAR data \cite{weng2021inverting, weng2022s2net, khurana2023point} to support decision-making in autonomous systems.

However, these research directions encounter substantial limitations in practical applications.
Existing works on LiDAR generation predominantly focus on learning the data distribution within a dataset, enabling them to synthesize only non-sequential, single LiDAR point clouds. 
In contrast, the autonomous driving field often requires temporally continuous LiDAR scene segments \cite{Geiger2013IJRR, caesar2020nuscenes}.
For LiDAR prediction tasks, with the rise of end-to-end autonomous driving (E2E-AD) \cite{hu2023planning, jiang2023vad, weng2024drive}, current methods are confronted with the following fundamental challenges:
\begin{enumerate}[label=\arabic*), leftmargin=*, align=left]
\item Prediction requires multiple frames of historical scene data; however, in closed-loop simulation \cite{chitta2022transfuser, dosovitskiy2017carla, jia2024bench2drive, prakash2021multi}, only the initial frame is available;
\item These methods are limited to predicting LiDAR scenes along a fixed trajectory, while the ego-vehicle may generate diverse trajectories based on different decisions;
\item They cannot accurately predicte LiDAR data over long temporal horizons, due to their inability to utilize diverse scene condition information in real time.
\end{enumerate} 

To address these challenges, we explore LiDAR-based generative world models and propose LaGen. 
LaGen is an autoregressive LiDAR point cloud generation framework capable of producing high-fidelity, long-horizon LiDAR scenes on a frame-by-frame basis.
Specifically, we first utilize spherical projection to convert the three-dimensional LiDAR data into a more compact form: a range image.
This representation not only improves computational efficiency but also better accommodates advanced diffusion models.
Next, we construct a high-fidelity LiDAR data generator based on the Latent Diffusion Model (LDM) \cite{rombach2022high}, and guide the diffusion process using multiple control conditions.

To further enhance the spatiotemporal consistency of the entire generative framework, we also introduce the following carefully designed strategies:
To improve the spatial consistency of generated results, we introduce a Scene Decoupling Estimation (SDE) module. 
It can provide the model with real-time object-level estimations of current scene, thereby enhancing the model's interactive generation capability.
To enhance the temporal consistency of generated results, we introduce a Noise Modulation (NM) module designed to alleviate error accumulation. 
This is because the autoregressive paradigm suffers from a train-inference mismatch; during inference, errors in the generated samples accumulate over longer temporal horizons.
The proposed module injects noise into the conditional features containing information from the previous frame, thereby reducing the model's over-reliance on prior frames and enabling it to better adapt to imperfect conditions.

We compare LaGen’s generator with several state-of-the-art LiDAR generation models on the nuScenes dataset. 
Experimental results demonstrate that our method achieves superior performance in generating high-fidelity LiDAR point clouds.
%In addition, we construct a protocol based on the nuScenes dataset for evaluating long-horizon LiDAR scene generation tasks.
In addition, we establish an evaluation protocol based on the nuScenes dataset for long-horizon LiDAR scene generation tasks.
Experimental results indicate that LaGen can accurately generate long-horizon LiDAR scenes using only a single-frame historical input.
It is noteworthy that our autoregressive framework further enables autonomous systems to incorporate decisions made at each time step into the generation of future predictions, thereby naturally supporting interactive world simulation.

In summary, our contributions are as follows:
\begin{itemize}[leftmargin=*]
\item We propose LaGen, which to the best of our knowledge is the first interactive framework capable of autoregressively generating high-fidelity, long-horizon LiDAR scenes;
\item We design a scene decoupling estimation module and a noise modulation module, which effectively enhance the model's performance in interactive scene-detail generation and long-horizon autoregressive generation;
\item We comprehensively evaluate the performance of the LaGen generator on the nuScenes dataset and establish a nuScenes-based benchmark for long-horizon generation tasks to demonstrate the superiority of our approach.
\end{itemize}